Throughout this section, $\excitation \geq 0$, $\baseIntensity \geq 0$ and
$\baseIntensity \neq 0$,
and $\branchingMatrix = \excitation/\decay$ as in \eqref{eq.GammaIsAOverBeta}.

Non-negativity here is entrywise: $\branchingMatrix \geq 0$ means
$\branchingMatrix\subscriptee \geq 0$ for every pair of types.
Neither $\excitation$ nor $\branchingMatrix$ is assumed symmetric.

\begin{lemma}[Perron-Frobenius]\label{lemma.perronFrobenius}
 Let $\branchingMatrix \geq 0$ be a square matrix with spectral radius $\branchingRatio$.
 Then
 \begin{enumerate}[label={\emph{\roman{*}}.}]
  \item $\branchingRatio$ is an eigenvalue of $\branchingMatrix$, and it is real
        and non-negative;
  \item there are non-negative, non-zero vectors $u$ and $v$ with
        $\branchingMatrix u = \branchingRatio u$ and
        $v^{\transpose}\branchingMatrix = \branchingRatio v^{\transpose}$,
        which we call a \emph{right} and a \emph{left Perron vector} of $\branchingMatrix$;
  \item $\branchingRatio$ lies between the smallest and the largest row sum of
        $\branchingMatrix$;
  \item if $\branchingRatio < 1$ then $I - \branchingMatrix$ is invertible and
        $(I-\branchingMatrix)^{-1} = \sum_{k\geq 0}\branchingMatrix^{k}$,
        a series with non-negative entries.
 \end{enumerate}
\end{lemma}
For a proof of \emph{i.}--\emph{iii.} see
\cite[Theorem 8.3.1 and Theorem 8.1.22]{HJ13mat};
the left Perron vector is the right Perron vector of $\branchingMatrix^{\transpose}$,
whose spectral radius is again $\branchingRatio$.
Under irreducibility $u$ and $v$ are strictly positive and $\branchingRatio$ is simple.
Part \emph{iv.} is the Neumann series, its entries non-negative because every term is.
Applying \emph{iii.} to $\branchingMatrix^{\transpose}$ bounds $\branchingRatio$ between the
smallest and the largest \emph{column} sum of $\branchingMatrix$;
in particular, constant column sums $c$ force $\branchingRatio = c$.

By Proposition \ref{prop.clusterRepresentation},
$\branchingMatrix\subscriptee$ is the expected number of direct offspring of type $e$
born of one event of type $\eone$.
The readable contraction is therefore the row vector
$\mathbf{1}^{\transpose}\branchingMatrix$,
whose $\eone$-th entry is the expected number of direct offspring of one event of type
$\eone$.

Recall that
a stochastic process $X$ 
is called 
\emph{stationary} 
when its law is invariant under time shift:
for every $r$ and every $t_1 < \dots < t_r$,
the joint law of $(X_{t_1+h}, \dots, X_{t_r+h})$ does not depend on $h$.

\begin{defi}\label{def.stationaryHawkes}
 A Hawkes process is \emph{stationary} if it admits a stationary version,
 that is, if there is a Hawkes process with the same baselines and kernel,
 defined on the whole line,
 whose law is invariant under time shifts.
\end{defi}

\begin{thm}[{\citealp{BM96sta}}]\label{thm.stationarity}
 Let $\multiCountingProc$ be a Hawkes process with integrable kernel
 $\hawkesKernel \geq 0$ and baselines $\baseIntensity > 0$ entrywise.
 Then the following are equivalent:
 \begin{enumerate}[label={\emph{\roman{*}}.}]
  \item $\multiCountingProc$ is stationary in the sense of
        Definition \ref{def.stationaryHawkes};
  \item $\branchingRatio < 1$.
 \end{enumerate}
 In that case the mean intensity is constant in time and
 \begin{equation}\label{eq.stationaryIntensity}
  \stationaryIntensity := \Expectation \big[ \intensity[ ](t) \big]
  = (I - \branchingMatrix)^{-1} \baseIntensity .
 \end{equation}
\end{thm}
For a proof, see \cite[Theorem 1]{BM96sta}.

The form of \eqref{eq.stationaryIntensity} follows from the definition.
A stationary version is defined on the whole line,
so \eqref{eq.intensity_ordHawkes} reads
$\intensity(t) = \baseIntensity_e + \sum_{\eone}\int_{-\infty}^{t}
\hawkesKernel\subscriptee(t-s)\,d\countingProc[\eone](s)$.
The kernel is non-negative, so Tonelli's theorem licenses the interchange of
expectation and integral, and stationarity makes
$\Expectation\big[d\countingProc[\eone](s)\big] = \stationaryIntensity[\eone]\,ds$,
whence
$\stationaryIntensity = \baseIntensity + \branchingMatrix\stationaryIntensity$.
By Lemma \ref{lemma.perronFrobenius}~\emph{iv.}, $\branchingRatio < 1$ makes
$I - \branchingMatrix$ invertible, so that equation has the unique solution
\eqref{eq.stationaryIntensity}, and the entries of the Neumann series are non-negative,
so $\stationaryIntensity \geq 0$.

\begin{remark}\label{remark.strictBaseline}
 Notice that $\baseIntensity > 0$ is what makes \emph{i.} imply \emph{ii.}
 Let $v$ be a left Perron vector of $\branchingMatrix$,
 non-negative and non-zero by Lemma \ref{lemma.perronFrobenius}~\emph{ii.}
 If the process is stationary with finite mean intensity, then
 $\stationaryIntensity = \baseIntensity + \branchingMatrix\stationaryIntensity$ ---
 the fixed point, which is available whatever $\branchingRatio$ is,
 where the inverse in \eqref{eq.stationaryIntensity} is not --- so that
 \begin{equation*}
  v^{\transpose}\baseIntensity
  = (1-\branchingRatio)\, v^{\transpose}\stationaryIntensity .
 \end{equation*}
 At $\branchingRatio \geq 1$ the right-hand side is non-positive
 and the left-hand side strictly positive, which is the contradiction.
 Without strict positivity the implication fails:
 $\branchingMatrix = \diag(2,0)$ with $\baseIntensity = (0,1)^{\transpose}$ has
 $\branchingRatio = 2$ and is stationary,
 because nothing reaches the first event type --- it has neither immigrants nor an exciter ---
 whereas
 the second component is a Poisson process.
\end{remark}

We write
\begin{equation*}
 \totalRate := \mathbf{1}^{\transpose} \stationaryIntensity,
 \qquad
 \bar{\baseIntensity} := \mathbf{1}^{\transpose} \baseIntensity
\end{equation*}
for the total event rate and the total immigration rate.

\begin{prop}\label{prop.meanState}
 Under the hypotheses of Theorem \ref{thm.stationarity},
 the stationary mean of the decayed counts is
 \begin{equation}\label{eq.meanState}
  m := \Expectation[\decayedCounts] = \frac{\stationaryIntensity}{\decay},
 \end{equation}
 and it holds
 \begin{equation}\label{eq.meanStateIdentity}
	 (  \decay I - \excitation) \, m = \baseIntensity .
 \end{equation}
\end{prop}
\begin{proof}
 Theorem \ref{thm.markovState} with $f(z) = z$ and \eqref{eq.intensityFromState} give
 $\frac{d}{dt}\Expectation[\decayedCounts(t)]
 = -\decay\Expectation[\decayedCounts(t)] + \Expectation[\intensity[ ](t)]
 = \baseIntensity + (\excitation - \decay I)\Expectation[\decayedCounts(t)]$.
 At stationarity the left-hand side vanishes, which is \eqref{eq.meanStateIdentity}.
\end{proof}

By Proposition \ref{prop.clusterRepresentation} the events partition into clusters,
one per immigrant.
What we measure next is the average size of a cluster,
and with it the share of the flow that the flow itself produced.

\begin{defi}\label{def.endogenousFraction}
 The \emph{endogenous fraction} of a stationary Hawkes process is
 \begin{equation}\label{eq.endogenousFraction}
  1 - \frac{\bar{\baseIntensity}}{\totalRate},
 \end{equation}
 the long-run share of events that are offspring rather than immigrants.
\end{defi}

\begin{prop}\label{prop.clusterSize}
 Assume $\branchingRatio < 1$.
 The expected size of the cluster founded by an immigrant of type $\eone$,
 counting the immigrant itself, is the $\eone$-th entry $C_{\eone}$ of
 \begin{equation}\label{eq.clusterSize}
  C = \mathbf{1}^{\transpose} (I - \branchingMatrix)^{-1} .
 \end{equation}
 The expected number of its strict descendants is $C_{\eone} - 1$.
 If moreover $\bar{\baseIntensity} > 0$, the mean cluster size is
 \begin{equation}\label{eq.meanClusterSize}
  \bar{C} := \sum_{\eone} \frac{\baseIntensity_{\eone}}{\bar{\baseIntensity}} C_{\eone}
  = \frac{\totalRate}{\bar{\baseIntensity}} ,
 \end{equation}
 and the endogenous fraction \eqref{eq.endogenousFraction} is $1 - 1/\bar C$.
\end{prop}
\begin{proof}
 The expected number of members of the $k$-th generation of the cluster, by type,
 is $\branchingMatrix^{k}\unitVector_{\eone}$: this holds for $k=0$, and each individual of
 type $\eone$ has on average $\branchingMatrix\subscriptee$ children of type $e$,
 independently of the rest, so the expectations propagate by one multiplication by
 $\branchingMatrix$ per generation.
 Summing over $k$ and contracting with $\mathbf{1}^{\transpose}$,
 all terms being non-negative so that the interchange is licit,
 gives $C_{\eone} = \mathbf{1}^{\transpose}\sum_{k\geq0}\branchingMatrix^{k}
 \unitVector_{\eone}$, which is \eqref{eq.clusterSize} by
 Lemma \ref{lemma.perronFrobenius}~\emph{iv.}
 The term $k=0$ is the immigrant.

 For the second half, immigrants of type $\eone$ arrive as a Poisson process of rate
 $\baseIntensity_{\eone}$, independently across $\eone$,
 so the type of a cluster's founder is distributed as
 $\baseIntensity_{\eone}/\bar{\baseIntensity}$.
 Then $\bar C = \mathbf{1}^{\transpose}(I-\branchingMatrix)^{-1}\baseIntensity /
 \bar{\baseIntensity}$, which is $\totalRate/\bar{\baseIntensity}$ by
 \eqref{eq.stationaryIntensity}.
\end{proof}

\begin{remark}\label{remark.clusterBudget}
 Notice that $\bar{\baseIntensity}\,\bar{C} = \totalRate$:
 immigrants arrive at rate $\bar{\baseIntensity}$,
 each founds a cluster of $\bar C$ events on average,
 and every event belongs to exactly one cluster.
 It is the cheapest available check that an implementation of
 \eqref{eq.stationaryIntensity} and \eqref{eq.clusterSize} is consistent.
\end{remark}

\begin{corol}\label{corol.constantClusterSize}
 Assume $\branchingRatio < 1$.
 The following are equivalent:
 \begin{enumerate}[label={\emph{\roman{*}}.}]
  \item $C_{\eone} = 1/(1-\branchingRatio)$ for every $\eone$;
  \item $\mathbf{1}^{\transpose}\branchingMatrix
        = \branchingRatio\mathbf{1}^{\transpose}$.
 \end{enumerate}
\end{corol}
\begin{proof}
 Since $\branchingRatio < 1$, the matrix $I - \branchingMatrix$ is invertible and
 $1 - \branchingRatio \neq 0$, so by \eqref{eq.clusterSize} the first statement says
 $\mathbf{1}^{\transpose}(I-\branchingMatrix)^{-1}
 = (1-\branchingRatio)^{-1}\mathbf{1}^{\transpose}$.
 Multiplying on the right by $I - \branchingMatrix$,
 this holds if and only if
 $(1-\branchingRatio)\mathbf{1}^{\transpose}
 = \mathbf{1}^{\transpose}(I - \branchingMatrix)$,
 which is the second statement.
\end{proof}

In general the branching ratio $\branchingRatio$
is neither the endogenous fraction nor $1-1/\bar C$.
Two conditions are sufficient for $\branchingRatio = 1 - 1/\bar C$, and neither is generic.
The first is $\mathbf{1}^{\transpose}\branchingMatrix
= \branchingRatio \mathbf{1}^{\transpose}$,
which by Corollary \ref{corol.constantClusterSize} also makes $C_{\eone}$ constant in
$\eone$.
The second is $\branchingMatrix\baseIntensity = \branchingRatio\baseIntensity$:
then $(I-\branchingMatrix)^{-1}\baseIntensity
= \baseIntensity/(1-\branchingRatio)$, so \eqref{eq.meanClusterSize} gives
$\bar C = 1/(1-\branchingRatio)$ with no condition on the column sums.
Notice that the second does not follow from the corollary,
and that it leaves $C_{\eone}$ free to vary with $\eone$.

\subsubsection{Second-order structure}
\label{sec.secondOrder}

What precedes describes the process in the mean.
We now compute its stationary second moments,
in closed form, for the exponential kernel.
Throughout, $\branchingRatio < 1$ and $\baseIntensity > 0$,
so that Theorem \ref{thm.stationarity} applies
and the process is taken in its stationary version.
The statements below are about that version;
the proofs reach it from $\decayedCounts(0) = 0$.

Under stationarity the covariance of $\decayedCounts(t)$ does not depend on $t$, and we write
\begin{equation}\label{eq.stationaryCovariance}
 V := \CoVariance\big(\decayedCounts(t)\big)
 = \Expectation\big[\decayedCounts(t)\decayedCounts^{\transpose}(t)\big]
 - m m^{\transpose},
\end{equation}
where $m$ is the mean decayed count of Proposition \ref{prop.meanState}.

Theorem \ref{thm.markovState} is applied with $f(z) = z z^{\transpose}$.
At an arrival of type $e$ the state gains the unit vector $\unitVector_e$, so that
\begin{equation}\label{eq.stateProductJump}
 \decayedCounts(\arrivalTimes_{j}+)\decayedCounts^{\transpose}(\arrivalTimes_{j}+)
 - \decayedCounts(\arrivalTimes_{j})\decayedCounts^{\transpose}(\arrivalTimes_{j})
 = \decayedCounts(\arrivalTimes_{j}) \unitVector_e^{\transpose}
 + \unitVector_e \decayedCounts^{\transpose}(\arrivalTimes_{j})
 + \unitVector_e \unitVector_e^{\transpose} .
\end{equation}
The jumps of $\decayedCounts$ are \emph{unit} vectors, which gives the last term
$\unitVector_e \unitVector_e^{\transpose}$;
and no two components jump together, by \eqref{eq.noCommonJumps},
which keeps that term diagonal.

\begin{defi}[Hurwitz matrix]\label{def.hurwitzMatrix}
 A real square matrix is \emph{Hurwitz} if all of its eigenvalues have strictly negative
 real part.
\end{defi}

\begin{prop}\label{prop.hurwitzCriterion}
 Let $\excitation \geq 0$ and $\decay > 0$.
 Then $\hurwitzMatrix := \excitation - \decay I$ is Hurwitz if and only if
 $\branchingRatio < 1$.
\end{prop}
\begin{proof}
 The eigenvalues of $\hurwitzMatrix$ are $\decay(\gamma_i - 1)$ over the eigenvalues
 $\gamma_i$ of $\branchingMatrix$,
 so $\hurwitzMatrix$ is Hurwitz if and only if $\Real \gamma_i < 1$ for every $i$.
 If $\branchingRatio < 1$, then
 $\Real \gamma_i \leq \lvert \gamma_i \rvert \leq \branchingRatio < 1$ for every $i$.
 Conversely $\branchingMatrix \geq 0$, so by
 Lemma \ref{lemma.perronFrobenius}~\emph{i.} the spectral radius $\branchingRatio$ is
 itself an eigenvalue of $\branchingMatrix$ and is real, whence $\branchingRatio < 1$.
\end{proof}

\begin{remark}\label{remark.hurwitzNeedsNonNegativity}
 Notice that $\excitation \geq 0$ is needed for the converse.
 Without it the implication runs one way only:
 $\branchingMatrix$ with eigenvalues $\pm 2i$ has $\branchingRatio = 2$
 while $\Real\gamma_i = 0 < 1$.
\end{remark}

\begin{thm}\label{thm.lyapunov}
 Let $\excitation \geq 0$, $\decay > 0$, $\baseIntensity > 0$ and $\branchingRatio < 1$,
 and write $\hurwitzMatrix := \excitation - \decay I$.
 Then the stationary version of the exponential-kernel Hawkes process has finite second
 moments, and its covariance $V$ of \eqref{eq.stationaryCovariance} is the unique solution
 of the Lyapunov equation
 \begin{equation}\label{eq.lyapunov}
  \hurwitzMatrix V + V \hurwitzMatrix^{\transpose} + \diag(\stationaryIntensity) = 0 .
 \end{equation}
 Moreover $V$ is positive definite and given by
 \begin{equation}\label{eq.explicitFormulaeOfStationaryCovariance}
  V = \int_{0}^{\infty} e^{\hurwitzMatrix s} \diag(\stationaryIntensity)
  e^{\hurwitzMatrix^{\transpose} s} \, ds .
 \end{equation}
 Finally, the stationary covariance of the intensity is
 $\CoVariance(\intensity[ ]) = \excitation V \excitation^{\transpose}$.
\end{thm}
\begin{proof}
 Start the process at $\decayedCounts(0) = \multiCountingProc(0) = 0$ and write
 $m_t := \Expectation[\decayedCounts(t)]$ and
 $\secondMoment_t := \Expectation[\decayedCounts(t)\decayedCounts(t)^{\transpose}]$,
 both finite for every $t$ by Remark \ref{remark.noSubcriticality}.
 Theorem \ref{thm.markovState} with $f(z)=z$ and \eqref{eq.intensityFromState} give
 $\dot m_t = \hurwitzMatrix m_t + \baseIntensity$,
 and with $f(z) = zz^{\transpose}$, using
 $\Expectation[\decayedCounts \intensity[ ]^{\transpose}]
 = m_t\baseIntensity^{\transpose} + \secondMoment_t\excitation^{\transpose}$,
 \begin{equation}\label{eq.secondMomentODE}
  \dot{\secondMoment}_t
  = \hurwitzMatrix \secondMoment_t + \secondMoment_t \hurwitzMatrix^{\transpose}
  + \baseIntensity m_t^{\transpose} + m_t \baseIntensity^{\transpose}
  + \diag\big( \baseIntensity + \excitation m_t \big) .
 \end{equation}
 By Proposition \ref{prop.hurwitzCriterion}, $\branchingRatio<1$ makes $\hurwitzMatrix$
 Hurwitz, so $m_t \to m = -\hurwitzMatrix^{-1}\baseIntensity
 = \stationaryIntensity/\decay$.
 Equation \eqref{eq.secondMomentODE} is then linear in $\secondMoment_t$ with an
 inhomogeneity that converges, and its homogeneous part
 $\secondMoment \mapsto \hurwitzMatrix\secondMoment
 + \secondMoment\hurwitzMatrix^{\transpose}$ is stable because $\hurwitzMatrix$ is
 Hurwitz, so $\secondMoment_t$ converges to the unique fixed point of the limiting
 equation, which exists and is unique by
 Lemma \ref{lemma.lyapunovUniqueness}~\emph{i.--ii.}
 A stationary version exists by Theorem \ref{thm.stationarity},
 and its second moment is that fixed point, the limit being independent of the initial state.
 Writing $\secondMoment = V + m m^{\transpose}$ and using
 $\hurwitzMatrix m = -\baseIntensity$,
 \begin{equation*}
  \hurwitzMatrix m m^{\transpose} + m m^{\transpose}\hurwitzMatrix^{\transpose}
  = -\baseIntensity m^{\transpose} - m\baseIntensity^{\transpose} ,
 \end{equation*}
 so the two cross terms of \eqref{eq.secondMomentODE} cancel and what remains is
 \eqref{eq.lyapunov}, with $\diag(\baseIntensity + \excitation m)
 = \diag(\stationaryIntensity)$ by \eqref{eq.intensityFromState}.
 The integral formula and positive definiteness are
 Lemma \ref{lemma.lyapunovUniqueness}~\emph{ii.},
 the forcing term being diagonal with strictly positive entries.
 Finally $\intensity[ ] = \baseIntensity + \excitation\decayedCounts$ with
 $\baseIntensity$ constant, so
 $\CoVariance(\intensity[ ]) = \excitation V \excitation^{\transpose}$.
\end{proof}

\begin{remark}
 Notice that the forcing term of \eqref{eq.lyapunov} is the diagonal matrix of stationary
 intensities:
 fluctuation is injected at the rate at which events occur,
 one unit of variance per event,
 and with no covariance between types because no two of them fire together.
 The homogeneous part transports it,
 and the equation is the balance between the two.
\end{remark}

\begin{lemma}\label{lemma.lyapunovUniqueness}
 Let $\hurwitzMatrix$ be a real $d \times d$ matrix.
 Then
 \begin{enumerate}[label={\emph{\roman{*}}.}]
  \item the map $V \mapsto \hurwitzMatrix V + V\hurwitzMatrix^{\transpose}$ is a bijection
        on symmetric matrices if and only if no two eigenvalues of $\hurwitzMatrix$ sum to
        zero. A Hurwitz $\hurwitzMatrix$ satisfies that criterion, every such sum
        having strictly negative real part;
  \item if $\hurwitzMatrix$ is Hurwitz then, for every symmetric $Q$, the equation
        $\hurwitzMatrix V + V\hurwitzMatrix^{\transpose} + Q = 0$ has the unique solution
        $V = \int_{0}^{\infty} e^{\hurwitzMatrix s} Q e^{\hurwitzMatrix^{\transpose} s}
        \, ds$, and $Q \succ 0$ implies $V \succ 0$.
 \end{enumerate}
\end{lemma}
\begin{proof}
 The eigenvalues of the map in \emph{i.} are the sums $\zeta_i + \zeta_j$ of the
 eigenvalues of $\hurwitzMatrix$, which gives the stated criterion;
 see \cite[Theorem 2.4.4.1]{HJ13mat}.
 For \emph{ii.}, the integral converges because $\hurwitzMatrix$ is Hurwitz,
 and differentiating $e^{\hurwitzMatrix s}Qe^{\hurwitzMatrix^{\transpose}s}$ and
 integrating over $[0,\infty)$ gives $-Q$ on the left-hand side of the equation.
\end{proof}

\subsubsection{The scalar case}

Contracting \eqref{eq.lyapunov} with $\mathbf{1}^{\transpose}$ on the left and
$\mathbf{1}$ on the right gives a scalar equation as soon as
$\mathbf{1}^{\transpose}\excitation$ is proportional to $\mathbf{1}^{\transpose}$,
that is, as soon as $\excitation$ has constant column sums.
Those sums are then necessarily $\decay\branchingRatio$:
by Lemma \ref{lemma.perronFrobenius}~\emph{iii.} applied to $\branchingMatrix^{\transpose}$,
constant column sums force $\branchingRatio$ to equal them.
Hence
$\mathbf{1}^{\transpose}\hurwitzMatrix
= -\decay(1-\branchingRatio)\mathbf{1}^{\transpose}$, and
\begin{equation}\label{eq.scalarVariance}
 \Variance \big( \mathbf{1}^{\transpose}\decayedCounts \big)
 = \frac{\totalRate}{2\decay(1-\branchingRatio)} .
\end{equation}
In one dimension this is
$\Variance(\decayedCounts) = \stationaryIntensity/(2\decay(1-\branchingRatio))$,
and it is the formula usually quoted.

\begin{remark}\label{remark.scalarVariance}
 Notice the scope of \eqref{eq.scalarVariance}.
 It concerns $\mathbf{1}^{\transpose}\decayedCounts$, the total decayed count,
 and not any single component;
 it is not a statement about the total intensity,
 which in the same special case has variance
 $\decay\branchingRatio^{2}\totalRate/(2(1-\branchingRatio))$;
 and it requires constant column sums, which are not generic.
 Neither parametrisation of Section \ref{sec.orderFlowModel} has them,
 and on the two of them \eqref{eq.scalarVariance} errs in \emph{opposite} directions:
 it understates $\mathbf{1}^{\transpose}V\mathbf{1}$ on the one whose column sums are
 further from constant, and overstates it on the one whose column sums are nearer.
 Neither the size of the discrepancy nor its sign can therefore be read off the column
 sums; both depend on the whole eigenstructure of $\excitation$ and on
 $\stationaryIntensity$.
 The discrepancy does not depend on $\decay$, since both sides scale as $1/\decay$.
\end{remark}

When the covariance is wanted, \eqref{eq.lyapunov} is a linear system in the
$\numEventTypes(\numEventTypes+1)/2$ distinct entries of $V$,
solved once.

\subsubsection{Timescales}

Everything so far has been stationary.
This subsection asks instead how the process moves towards stationarity:
what the mean of the state at a future time is, given the state now,
and how long it takes the memory of that state to decay.
Section \ref{sec.priceFormation} reads the price off the first of these.

\begin{prop}\label{prop.forwardMean}
 Let $\multiCountingProc$ be an exponential-kernel Hawkes process,
 and write $m_{t}(z) := \Expectation_{z}\big[\decayedCounts(t)\big]$
 for the mean state at time $t$ of the process started from $\decayedCounts(0+) = z$
 and $\multiCountingProc(0) = 0$.
 Then
 \begin{equation}\label{eq.forwardMeanState}
  m_{t}(z) = e^{\hurwitzMatrix t} z
  + \Big( \int_{0}^{t} e^{\hurwitzMatrix s} \, ds \Big) \baseIntensity ,
 \end{equation}
 and, writing $\meanResponseIntegral(\forecastHorizon)
 := \int_{0}^{\forecastHorizon} e^{\hurwitzMatrix s} \, ds$,
 \begin{equation}\label{eq.forwardMeanCounts}
  \Expectation_{z}\big[ \multiCountingProc(\forecastHorizon) \big]
  = \forecastHorizon \, \baseIntensity
  + \excitation \, \meanResponseIntegral(\forecastHorizon) \, z
  + \excitation \Big( \int_{0}^{\forecastHorizon} \meanResponseIntegral(s) \, ds \Big)
  \baseIntensity .
 \end{equation}
 If $\branchingRatio < 1$, then
 $\meanResponseIntegral(\forecastHorizon)
 = \hurwitzMatrix^{-1}\big( e^{\hurwitzMatrix \forecastHorizon} - I \big)$.
\end{prop}
\begin{proof}
 Both means are finite at every $\branchingRatio$ and every finite horizon, by
 Remark \ref{remark.noSubcriticality}.
 Theorem \ref{thm.markovState} with $f(z) = z$, together with
 \eqref{eq.intensityFromState}, gives
 $\dot m_{t} = \hurwitzMatrix m_{t} + \baseIntensity$ with $m_{0+} = z$,
 and variation of constants gives \eqref{eq.forwardMeanState}.
 The compensator of $\multiCountingProc$ integrates the intensity, so
 $\Expectation_{z}[\multiCountingProc(\forecastHorizon)]
 = \int_{0}^{\forecastHorizon}
 \big( \baseIntensity + \excitation m_{s}(z) \big) \, ds$,
 which is \eqref{eq.forwardMeanCounts}.
 The closed form of $\meanResponseIntegral$ needs $\hurwitzMatrix$ invertible,
 which Proposition \ref{prop.hurwitzCriterion} supplies.
\end{proof}

\begin{remark}\label{remark.transientAgainstStationary}
 Notice that $m_{t}(z)$ is the \emph{transient} mean and is not the stationary mean $m$ of
 Proposition \ref{prop.meanState}, which does not depend on $t$.
 The two meet at the fixed point of the same equation:
 $\dot m_{t} = 0$ exactly at
 $m = -\hurwitzMatrix^{-1}\baseIntensity = \stationaryIntensity/\decay$,
 and when $\branchingRatio < 1$ the matrix $\hurwitzMatrix$ is Hurwitz, so $m_t(z)$
 converges to $m$ from every $z$.
\end{remark}

\begin{corol}\label{corol.meanResponse}
 Let $\branchingRatio < 1$.
 Consider two copies of the process differing at time $0$ only in that the state of the
 second carries one extra event of type $\eone$, and write
 $R_{\eone}(t)$ for the difference of their mean states at time $t$.
 Then
 \begin{equation}\label{eq.meanResponse}
  R_{\eone}(t) = e^{\hurwitzMatrix t}\unitVector_{\eone} ,
  \qquad
  \mathbf{1}^{\transpose}R_{\eone}(t)
  = \mathbf{1}^{\transpose} e^{\hurwitzMatrix t}\unitVector_{\eone} ,
 \end{equation}
 a combination of exponentials at the rates $\decay(1 - \Real\gamma_i)$ when
 $\branchingMatrix$ is diagonalisable, $\gamma_i$ running over its eigenvalues,
 of which the slowest rate is $\decay(1-\branchingRatio)$.
 In general the modes carry polynomial factors as well.
 The corresponding response of the mean intensity is
 $\excitation e^{\hurwitzMatrix t}\unitVector_{\eone}$, and
 \begin{equation}\label{eq.responseIntegral}
  \int_{0}^{\infty} \mathbf{1}^{\transpose}\excitation e^{\hurwitzMatrix t}
  \unitVector_{\eone} \, dt
  = \mathbf{1}^{\transpose}\big[ (I - \branchingMatrix)^{-1} - I \big] \unitVector_{\eone}
  = C_{\eone} - 1 .
 \end{equation}
\end{corol}
\begin{proof}
 By \eqref{eq.forwardMeanState} the two mean states are
 $m_{t}(z)$ and $m_{t}(z + \unitVector_{\eone})$,
 whose difference is $e^{\hurwitzMatrix t}\unitVector_{\eone}$:
 the inhomogeneity is the same for both copies and cancels.
 The eigenvalues of $\hurwitzMatrix$ are $\decay(\gamma_i - 1)$, giving the stated rates;
 by Lemma \ref{lemma.perronFrobenius}~\emph{i.} the spectral radius $\branchingRatio$ is
 itself an eigenvalue of $\branchingMatrix$ and dominates the others in modulus,
 so $\decay(1-\branchingRatio)$ is the smallest of them.
 The intensity response follows from
 $\intensity[ ] = \baseIntensity + \excitation\decayedCounts$ with $\baseIntensity$
 constant.
 For \eqref{eq.responseIntegral}, $\meanResponseIntegral(\infty) = -\hurwitzMatrix^{-1}
 = \decay^{-1}(I-\branchingMatrix)^{-1}$,
 and $\excitation\decay^{-1} = \branchingMatrix$ gives
 $\branchingMatrix(I-\branchingMatrix)^{-1} = (I-\branchingMatrix)^{-1} - I$;
 contracting with $\mathbf{1}^{\transpose}$ and using \eqref{eq.clusterSize} gives
 $C_{\eone}-1$.
\end{proof}

\begin{defi}\label{def.relaxationTime}
 The \emph{relaxation time} of a subcritical Hawkes process is
 \begin{equation}\label{eq.relaxationTime}
  \frac{1}{\decay(1-\branchingRatio)} .
 \end{equation}
\end{defi}


\begin{remark}
 Notice that \eqref{eq.relaxationTime} is the time constant of the \emph{slowest} mode of
 \eqref{eq.meanResponse}, and describes the response only asymptotically.
 It is not a half-life; the half-life of the slowest mode is
 $\log 2 / (\decay(1-\branchingRatio))$.
 Nor need the response be monotone: by \eqref{eq.meanResponse},
 $\frac{d}{dt}\big\vert_{t=0}\mathbf{1}^{\transpose}R_j(t)
 = \decay\big( (\mathbf{1}^{\transpose}\branchingMatrix)_j - 1 \big)$,
 so the response to an event of type $j$ rises before it falls precisely when the $j$-th
 column sum of $\branchingMatrix$ exceeds one.
 For \eqref{eq.meanResponse} to be asymptotically \emph{proportional} to
 $e^{-\decay(1-\branchingRatio)t}$, with a strictly positive constant for every $j$,
 the Perron mode must appear in every response, which holds when $\branchingMatrix$ is
 irreducible; otherwise $\decay(1-\branchingRatio)$ only bounds the decay.
\end{remark}

\begin{defi}\label{def.dimensionlessGroups}
 For a stationary Hawkes process of total rate $\totalRate$ observed in windows of
 length $w$, we call $\totalRate/\decay$ the \emph{clustering number}
 and $\totalRate w$ the \emph{window count}.
\end{defi}

\begin{remark}\label{remark.dimensionlessGroups}
 Notice that these two numbers, and not $\decay$ or $\totalRate$ separately, are what a
 measurement made on window sums depends on.
 The clustering number is the number of events arriving within one excitation timescale,
 so it decides whether self-excitation is visible at the rate the process is sampled:
 at $\totalRate/\decay \ll 1$ the excitation of one event has expired before the next
 arrives, and a study conducted there measures a process in which there is correctly
 nothing to find.
 The window count is the number of events falling in one window,
 and it decides whether a window sum is a statistic or a single event.
 Both are invariant under the rescaling $t \mapsto ct$,
 $(\baseIntensity,\excitation,\decay,w) \mapsto
 (\baseIntensity/c, \excitation/c, \decay/c, cw)$,
 which leaves $\branchingMatrix$ and hence $\branchingRatio$ unchanged.
\end{remark}

\subsubsection{Sub- and supercriticality}

Everything above assumes $\branchingRatio < 1$.
This subsection separates what that assumption is needed for from what holds without it.

Three things hold at every $\branchingRatio$.
For a linear Hawkes process and every finite horizon $\timeHorizon$,
$\Expectation[\countingProc(\timeHorizon)^2] < \infty$
by the generation bound of Remark \ref{remark.noSubcriticality};
the forward means \eqref{eq.forwardMeanState} and \eqref{eq.forwardMeanCounts} are finite,
for the same reason;
and on $[0,\timeHorizon]$ the process
$\martingale := \multiCountingProc - \compensator$
is a square-integrable martingale,
with $\langle \martingale_e \rangle = \compensator_e$
and $\langle \martingale_e, \martingale_{\eone} \rangle \equiv 0$ for $e \neq \eone$,
the latter by \eqref{eq.noCommonJumps}.
Subcriticality does not enter any of the three.

Three others need it, and each fails without it.
A stationary version exists, by Theorem \ref{thm.stationarity};
the second moments of the state $(\decayedCounts, \intensity[ ])$ are bounded uniformly in
$\timeHorizon$, by Theorem \ref{thm.lyapunov};
and the process is ergodic, so that time averages along one path estimate the stationary
quantities of this section.
Each fails at $\branchingRatio \geq 1$, provided $\baseIntensity > 0$ ---
Remark \ref{remark.strictBaseline} is the counterexample without it.
Contracting $\dot m_{t} = \hurwitzMatrix m_{t} + \baseIntensity$ with a left Perron
vector $v$, and using $v^{\transpose}\excitation = \decay\branchingRatio\, v^{\transpose}$,
\begin{equation}\label{eq.supercriticalGrowth}
 v^{\transpose} \Expectation \big[ \intensity[ ](t) \big]
 = \begin{cases}
   \big( v^{\transpose}\baseIntensity \big)\,(1 + \decay t),
   & \branchingRatio = 1, \\[2pt]
   v^{\transpose}\baseIntensity \,
   \dfrac{\branchingRatio\, e^{\decay(\branchingRatio-1)t} - 1}{\branchingRatio - 1},
   & \branchingRatio > 1 .
 \end{cases}
\end{equation}
Above criticality the intensity grows exponentially at rate $\decay(\branchingRatio-1)$.
At $\branchingRatio = 1$ it grows at least linearly,
and linearly with a quadratic event count when $1$ is a semisimple eigenvalue of
$\branchingMatrix$, in particular when $\branchingMatrix$ is irreducible.

\begin{remark}\label{remark.martingaleNotBounded}
 Notice that the martingale is not bounded in $L^{2}$ on $[0,\infty)$ even when
 $\branchingRatio<1$:
 $\Expectation[\martingale_e(\timeHorizon)^2]
 = \Expectation[\compensator_e(\timeHorizon)]
 \sim \stationaryIntensity[e] \timeHorizon$ diverges.
 What subcriticality bounds is the state, not the martingale.
\end{remark}
